%% Chapter 3 — Background
\chapter{Background}
\label{ch:background}

This chapter introduces the technical and theoretical foundations required for the work described in subsequent chapters. It covers the concept of virtual avatars and visual representation, edge detection mathematics, the Kuwahara painterly filter, and the human factors concepts of trust and social presence relevant to the user study.

\section{Mathematical Notation}
\label{sec:notation}

Throughout this thesis, scalars are denoted by italic letters~$x$, vectors by bold lowercase letters~$\mathbf{x}$, and matrices by bold uppercase letters~$\mathbf{X}$. The $i$-th element of vector $\mathbf{x}$ is written $x_i$. Convolution of a signal $I$ with a kernel $K$ is written $K * I$. In shader pseudocode, \texttt{ddx(f)} and \texttt{ddy(f)} refer to the hardware-computed screen-space partial derivatives of a value~\texttt{f} with respect to the horizontal and vertical screen axes, respectively.

\section{Edge Detection}
\label{sec:background_edge}

An edge in an image is conventionally defined as a location where intensity changes rapidly. Detecting such locations involves estimating the spatial gradient of the image intensity function, then applying a threshold to identify pixels where the gradient magnitude exceeds some criterion.

\subsection{Gradient-Based Detection}
\label{sec:background_gradient}

Given an image $I(x, y)$, the gradient is the vector $\nabla I = (\partial I / \partial x, \, \partial I / \partial y)$. Its magnitude
\begin{equation}
  M = \|\nabla I\| = \sqrt{\left(\frac{\partial I}{\partial x}\right)^2 + \left(\frac{\partial I}{\partial y}\right)^2}
  \label{eq:gradient_magnitude}
\end{equation}
provides a scalar measure of edge strength at each pixel. In practice, the partial derivatives are approximated by discrete convolution with finite-difference kernels.

The Sobel operator~\cite{gonzalez2018} estimates the gradient using two $3\times3$ kernels:
\begin{equation}
  G_x = \begin{bmatrix} -1 & 0 & +1 \\ -2 & 0 & +2 \\ -1 & 0 & +1 \end{bmatrix}, \qquad
  G_y = \begin{bmatrix} +1 & +2 & +1 \\ 0 & 0 & 0 \\ -1 & -2 & -1 \end{bmatrix}.
  \label{eq:sobel_kernels}
\end{equation}
$G_x$ responds to horizontal intensity transitions (vertical edges); $G_y$ responds to vertical transitions. The gradient magnitude is $M = \sqrt{G_x^2 + G_y^2}$, computed pointwise after convolving the image with each kernel separately. The $\pm2$ weights along the central row and column give the Sobel operator a small amount of spatial smoothing relative to the simpler Prewitt kernel, improving noise robustness marginally.

The Roberts Cross operator uses a $2\times2$ diagonal kernel pattern instead:
\begin{equation}
  R_1 = \begin{bmatrix} +1 & 0 \\ 0 & -1 \end{bmatrix}, \qquad
  R_2 = \begin{bmatrix} 0 & +1 \\ -1 & 0 \end{bmatrix}.
  \label{eq:roberts_kernels}
\end{equation}
The edge magnitude is $\lvert R_1 \rvert + \lvert R_2 \rvert$, which approximates the gradient diagonal without a centre sample. Roberts Cross is faster and more local than Sobel but more sensitive to noise. In this thesis it is used specifically for colour edge detection within the hierarchical technique, where its two-sample footprint keeps the per-fragment sample count low.

\subsection{Screen-Space Normal Derivatives}
\label{sec:background_normal_deriv}

Modern GPU pipelines expose per-fragment screen-space derivatives of any interpolated value via the \texttt{ddx} and \texttt{ddy} intrinsics, computed using the finite difference between adjacent $2\times2$ pixel quads. Applying this to the world-space surface normal $\mathbf{N}$ gives
\begin{equation}
  M_N = \sqrt{\|\texttt{ddx}(\mathbf{N})\|^2 + \|\texttt{ddy}(\mathbf{N})\|^2},
  \label{eq:normal_gradient}
\end{equation}
which responds to abrupt changes in surface orientation — creases, mesh boundaries, and regions of high curvature. This computation requires no additional texture samples and executes at negligible cost relative to the surrounding fragment shader. The significant practical limitation is that screen-space derivatives are measured in pixel units. A fixed-world-space edge — the crease between a jacket lapel and its collar, for instance — projects to a varying number of pixels depending on camera distance, causing edge thickness to scale with the viewpoint rather than with the surface.

\subsection{Gaussian Pre-Filtering}
\label{sec:background_gaussian}

Differentiation amplifies high-frequency components. Marr and Hildreth~\cite{marr1980} established that stable edge detection requires low-pass smoothing before gradient computation. Because convolution and differentiation are both linear operations, the smoothed gradient can be written as
\begin{equation}
  \nabla(G_\sigma * I) = (\nabla G_\sigma) * I,
  \label{eq:gauss_deriv_commute}
\end{equation}
where $G_\sigma$ is a Gaussian kernel with standard deviation $\sigma$. Canny's influential edge detector~\cite{canny1986} built on this principle: smooth with a Gaussian, compute the gradient of the smoothed image, and apply non-maximum suppression to obtain thin, stable edge maps. In a single-pass fragment shader, a true two-pass separable convolution is not available. The Gaussian can instead be approximated by sampling the source texture at a fixed set of offsets with precomputed weights, accepting that the smoothed values are recomputed independently at each sample position.

\section{The Kuwahara Filter}
\label{sec:background_kuwahara}

The Kuwahara filter~\cite{kyprianidis2009} is a structure-preserving smoothing operator that produces piecewise-uniform regions with sharp inter-region boundaries — a visual effect that resembles painterly brushwork. It operates by dividing the neighbourhood around each pixel into $k$ overlapping subregions, computing the mean and variance of pixel values within each subregion, and setting the output to the mean of whichever subregion has the lowest variance. Formally, for subregions $\{Q_i\}_{i=1}^{k}$ and their respective means $\mu_i$ and variances $\sigma_i^2$:
\begin{equation}
  F_{\text{Kuwahara}}(x, y) = \mu_j, \qquad j = \arg\min_i \, \sigma_i^2.
  \label{eq:kuwahara}
\end{equation}

The effect is that pixels on uniform surfaces are assigned the mean colour of their most locally uniform neighbourhood, while pixels near strong boundaries pick the subregion that falls entirely on one side. Edges are therefore sharpened rather than blurred, even as internal regions are smoothed toward a uniform painterly fill.

The classical isotropic formulation uses four quadrant-shaped subregions of equal size, producing a rotationally symmetric filter that does not prefer any orientation. An anisotropic extension by Kyprianidis et al.~\cite{kyprianidis2009} aligns the filter kernel with local image structure by computing a structure tensor from a prior render pass, producing brush-stroke effects oriented along salient contours. This produces superior results visually, but requires an intermediate render target that is incompatible with a single-pass fragment hook.

\section{Virtual Avatars and Visual Representation}
\label{sec:background_avatars}

In virtual environments, an avatar is a digital body that represents a user within a shared space. The term covers a wide range of forms, from abstract geometric markers to fully articulated, photorealistic human figures, and the choice of representation carries consequences that extend well beyond visual preference.

Avatar representations can be placed along a continuum of anthropomorphism and visual fidelity. At one end, abstract or stylised avatars retain the structural properties of a humanoid figure without committing to the surface detail of a real person. At the other, photorealistic avatars attempt to reproduce the user's physical appearance with enough accuracy to approach the visual fidelity of a photograph. Where an avatar sits along this continuum has measurable consequences for how users are perceived, how much social presence is felt in the interaction, and how the avatar is trusted~\cite{latoschik2017, nowak2003}.

Avatar appearance also shapes the behaviour of the person embodying it. Yee and Bailenson~\cite{yee2007} demonstrated through a series of experiments that users assigned more attractive or physically taller avatars adopted behaviours consistent with those characteristics: they disclosed more personal information, negotiated more assertively, and maintained shorter interpersonal distances. This tendency, termed the Proteus effect, illustrates that visual self-representation in virtual environments is not a passive interface element but an active influence on how users act and present themselves.

Contemporary avatar platforms have made high-quality avatar creation accessible without requiring manual artistic work. Services such as Avaturn reconstruct a personalised, fully rigged avatar from a single uploaded photograph using learned three-dimensional reconstruction, producing a mesh with facial blendshape animation support that integrates directly into standard game engine pipelines. This accessibility makes individualized photorealistic representation practical for a wide range of applications, and it sharpens the perceptual stakes: an avatar that looks like a real person will inevitably be compared to one, and any shortfall in that comparison has consequences both for immersion and for how the user is received by others.

\section{Trust and Social Presence in VR}
\label{sec:background_trust}

Trust in the context of human-avatar interaction is typically understood as a social judgement — the degree to which an observer is willing to rely on the entity represented by the avatar. It is related to but distinct from social presence, which refers to the subjective sense of sharing a space with another being. Both constructs are influenced by avatar appearance, but through different mechanisms. Social presence depends strongly on the perceived anthropomorphism of the avatar and the responsiveness of its motion~\cite{nowak2003}. Trust is additionally sensitive to perceived competence, predictability, and familiarity — including the absence of uncanny valley discomfort~\cite{mori1970}.

In user studies, social presence is most commonly measured with self-report instruments including the Social Presence Questionnaire~\cite{nowak2003} and the items from Witmer and Singer's sense-of-presence measure. Trust has been assessed with scales drawn from social psychology adapted for HCI contexts, where items probe dimensions such as reliability, benevolence, and integrity. The Godspeed questionnaire~\cite{macdorman2006} provides a validated set of semantic differential items for rating human-likeness, animacy, and eeriness, which have been used to characterise uncanny valley responses to robotic and avatar stimuli.

A consistent finding across this literature is that the relationship between visual fidelity and perceived trustworthiness is not linear. Near-photorealistic representations can produce lower trust scores than clearly stylised ones, consistent with the uncanny valley hypothesis. This non-linearity motivates the experimental design of this thesis, which compares a range of NPR abstraction levels rather than simply contrasting a stylised condition against the photorealistic baseline.
